% !TEX root = ../main.tex

\begin{acknowledgements}

  行文至此，本科阶段的学业已近尾声。本科四年，我笑过、哭过，迷茫过、挣扎过，也坚定过。本科的最后半年，一切都在变好，科研上有了继续深耕的环境，生活也看见了未来的希望。
  
  感谢导师在研究方向选择、论文写作和实验设计中给予的悉心指导与持续帮助。

  感谢 NoMaD 项目团队开源的 visualnav-transformer 代码仓库，为本文的研究提供了坚实的算法基础。

  感谢云深处科技提供的 Lite3 四足机器人平台与技术支持，使本文能够将视觉导航研究从仿真推进到真机验证。

  感谢实验室师兄在日常讨论、代码调试与真机实验中给予的无私帮助。初入科研时，我曾倍感忐忑，师兄们耐心指点使我突破瓶颈，顺利完成本课题。

  感谢上海交通大学自动化与感知学院提供的学习环境与科研条件。

  感谢智能感知工程专业，作为专业首届本科生，我曾无数次地怀疑这个专业是否能够给本科生足够的培养，但四年的学习与实践让我见证了它的快速成长与扎实积淀。在此，我不仅收获了系统的专业知识与工程能力，更结识了良师益友。

  感谢自感2253班级的兄弟姐妹们，感谢2253卷帘门。每逢期末的相互扶持与挑灯夜战，必将成为我青春岁月中最温暖的印记。

  感谢挚友同行，四载交大岁月，因你们而熠熠生辉。
  
  感谢家人一直以来的理解与支持。感谢爸爸妈妈，感谢你们二十多年的养育之恩，以及始终如一的理解与鼓励，是我前行路上最坚实的底气。

  感谢我的女朋友，这半年来，在我最困顿迷茫之时，是你的陪伴与鼓励让我重拾方向；无数个疲惫的夜晚，因你的理解与支持而不再难熬。你的温柔与坚定，是我坚持至今的重要力量。

  最后感谢一直走在路上的我自己，我没有被困难打倒，我没有被压力击溃，我一直在坚持，我一直在成长。我是好样的！愿我带着这份坚韧与热爱，继续奔赴下一程山海！
  
\end{acknowledgements}
